% ====================================================================
%  CHAPTER: CONCLUSIONS AND FUTURE WORK
% ====================================================================
\chapter{Conclusions}
\label{ch:conclusions}

% ====================================================================
%  §1  SUMMARY
% ====================================================================
\section{Summary}
\label{sec:concl-summary}

This thesis studied the impact of layer fusion on spatial DNN
accelerators by extending the FactorFlow analytical mapping framework
with support for fused multi-layer execution.
The work encompassed modeling, implementation, validation, and an
extensive experimental evaluation covering two spatial architectures
(DepFiN-like and Eyeriss-v1-like), four CNN workloads (FSRCNN,
MC-CNN, VGG16, ResNet18), three fusion modes (full, partial, single-layer),
and two architecture-sizing scenarios (Scenario~A, full-fusion-sized;
Scenario~B, partial-fusion-sized).
This section recapitulates the key findings along two axes: the
impact of layer fusion as an execution strategy, and the role that
the extended FactorFlow model plays in enabling its analysis.

% ------------------------------------------------------------------
\subsection{Impact of Layer Fusion}
\label{sec:concl-fusion-impact}

The central premise of layer fusion is that eliminating the
intermediate-activation round-trips through external memory can
reduce both energy and latency.
The experiments confirm part of this premise and, equally
importantly, reveal its boundaries: the outcome depends on the
workload's activation-to-weight traffic ratio, the architecture's
memory organization, and the per-access energy cost of DRAM.

\paragraph{DRAM traffic reduction.}
Full fusion achieves its primary objective: DRAM traffic is reduced
by $25$--$95\%$ depending on the workload ($94.8\%$ for FSRCNN,
$85.3\%$ for MC-CNN, $59.9\%$ for VGG16, $25.1\%$ for ResNet18).
Activation-dominant workloads benefit the most because their
intermediate activations constitute the overwhelming majority of
off-chip traffic in single-layer mode.
These reductions are identical across the two architectures because
the DRAM access pattern is determined by the tiling of the mapping,
and we use the same tile sizes on both architectures for fairness
of evaluation.

\paragraph{Latency benefit.}
The DRAM-traffic reduction translates into substantial latency
improvements for activation-dominant workloads: $42$--$48\%$ on
DepFiN and $16$--$46\%$ on Eyeriss.
For weight-dominant workloads the results depend on the
architecture.
On Eyeriss, full fusion still reduces latency ($7.8$--$21.2\%$),
because the distributed per-PE register architecture avoids a
single shared-memory contention point for all layers' weights.
On DepFiN, however, full fusion is \emph{slower} than single-layer
execution (up to $30.3\%$ latency increase for ResNet18) because
the very large shared weight memory (10.7--14.4\,MB) creates a
bandwidth bottleneck that outweighs the DRAM savings.
Partial (pairwise) fusion captures $87$--$103\%$ of the full-fusion
latency benefit, indicating that nearly all the inter-layer traffic
saving is realized by fusing just two consecutive layers.

\paragraph{Energy: a workload-dependent and architecture-dependent outcome.}
Unlike the simpler narrative of ``fusion always increases energy,''
the outcome depends critically on the interplay between the DRAM
energy saved and the on-chip memory overhead incurred, and
consequently on the DRAM per-access energy cost assumed in the
model.
At the 160\,pJ\,/\,byte DRAM cost used in this work
(Section~\ref{sec:energy-model}), the results split along workload
type and architecture:

\begin{itemize}
  \item On \textbf{DepFiN}, full fusion is energy-beneficial for
        activation-dominant workloads: FSRCNN single-layer energy is
        $4.38\times$ the full-fusion value, and MC-CNN is
        $1.73\times$.
        The eliminated DRAM traffic is large (576\,KB and 522\,KB
        feature memories are modest), so the on-chip overhead is
        outweighed by the DRAM savings.
  \item For weight-dominant workloads on DepFiN, the picture
        reverses: VGG16 single-layer energy is only $7.4\%$ of the
        full-fusion value, and ResNet18 is $17.3\%$.
        Fusion requires weight memories of 14.4\,MB (VGG16) or
        10.7\,MB (ResNet18), whose per-access energy compounds
        across billions of MACs, far exceeding the DRAM savings.
  \item On \textbf{Eyeriss}, full fusion is roughly energy-neutral
        for FSRCNN ($0.985\times$\,full) and energy-adverse for all
        other workloads (MC-CNN single costs $45\%$ of full, VGG16
        costs $5.9\%$, ResNet18 costs $12.6\%$).
        The replicated per-PE register architecture amplifies the
        on-chip energy: with 16\,384~PEs for VGG16, the aggregate
        weight register capacity is $1{,}200 \times 16{,}384 =
        19.7$\,M~bytes, each register access contributing to
        per-access energy.
\end{itemize}

\noindent
A lower DRAM energy cost would reduce the benefit of eliminating
DRAM traffic and shift the balance further against fusion; a higher
DRAM cost would extend fusion's energy advantage to more
workload--architecture combinations.
The success of fusion is therefore sensitive to the DRAM technology
node and packaging: designs with high-bandwidth memory (HBM) or
tightly integrated DRAM may find the energy case for fusion weaker
than those relying on conventional off-chip DRAM interfaces.

\paragraph{Energy--Delay Product.}
When energy and latency are combined into the EDP metric, full
fusion wins in three of the eight architecture--workload
combinations: FSRCNN on DepFiN ($6.44\times$ single/full), MC-CNN
on DepFiN ($2.31\times$), and FSRCNN on Eyeriss ($1.75\times$).
In all three, the activation-dominant nature of the workload
ensures that both energy and latency favor fusion, amplifying the
EDP advantage.
For the remaining five combinations, single-layer execution is the
more efficient choice, with the weight-dominant workloads showing
particularly unfavorable ratios (VGG16 on Eyeriss:
$0.074\times$; ResNet18 on DepFiN: $0.073\times$).

This result clarifies the role of layer fusion:
\emph{fusion is primarily a latency optimization whose energy
trade-off is workload-dependent}.
Its benefit grows with the activation-to-weight traffic ratio and
diminishes as the on-chip memory overhead increases.
Whether the trade-off is acceptable depends on the design priority,
the workload class, and the external memory technology:
throughput-critical deployments with activation-heavy workloads may
prefer the latency and EDP improvement, while energy-constrained
systems or weight-dominant workloads should favor single-layer
execution.

\paragraph{Partial fusion as a practical compromise.}
Partial fusion consistently captures most of the latency benefit
while reducing the memory and energy cost.
In Scenario~B, sizing the architecture for pairwise rather than
full fusion reduced the weight memory by up to $3.1\times$ on DepFiN
(VGG16: 14\,366\,KB${\to}$4\,610\,KB; ResNet18:
10\,738\,KB${\to}$4\,608\,KB) and the per-PE WReg on Eyeriss
(VGG16: 1\,200${\to}$576; ResNet18: 902${\to}$384).
The latency benefit was almost fully preserved: partial fusion
captured $96$--$100\%$ of the full-fusion latency saving for
activation-dominant workloads, and it eliminated or greatly reduced
the worst-case latency penalty on DepFiN
($-11.3\% \to +0.8\%$ for VGG16; $-30.3\% \to -7.5\%$ for
ResNet18).
The EDP verdict was identical in scope: three of eight
architecture--workload combinations favored fusion in both
scenarios.
For activation-dominant workloads on DepFiN, partial fusion
won on energy, latency, and EDP simultaneously, making it the
strictly dominant strategy.

\paragraph{Super-linear scaling of the energy penalty.}
The fixed-config scaling experiment on Eyeriss
(Section~\ref{sec:fusion-fixed-scaling}) revealed that the
fusion-to-single energy ratio is not a fixed constant: it
\emph{scales super-linearly} with register size.
Each $2\times$ increase in per-PE WReg raised the ratio by
$1.5$--$1.9\times$ (ResNet18: $8.0\times \to 11.7\times \to
19.6\times$; VGG16: $17.0\times \to 28.9\times \to 54.1\times$).
The growth arises because the fused energy increases with WReg
(each access touches a larger SRAM structure) while the
single-layer energy decreases only modestly (the DRAM component is
PE-count-independent).
A designer who shrinks the PE array to save area will face a
proportionally larger WReg and, consequently, a worse
fusion-to-single energy ratio.
This scaling behavior reinforces the workload-dependent nature of
the fusion trade-off and underlines the importance of accurate
energy models for both on-chip and off-chip memory when evaluating
fusion strategies.

% ------------------------------------------------------------------
\subsection{Impact of Using FactorFlow}
\label{sec:concl-factorflow-impact}

The choice to extend FactorFlow, rather than building a new tool or
adapting Timeloop, proved instrumental in several respects.

\paragraph{Generality of the factor-based abstraction.}
FactorFlow's representation of the map space as a set of factors
distributed across hierarchy levels and dimensions generalized
naturally to the fused-layer setting.
The aliasing mechanism (shared activation dimensions with concrete
intermediate names) and the decoupling mechanism (per-layer
disjoint dimension sets) were expressed entirely within the existing
factor vocabulary, without requiring changes to the core data
structures.

\paragraph{Map-space characterization.}
The aliasing and decoupling concepts enabled a precise
characterization of how the fused map space grows with the number
of layers.
Each additional fused layer introduces six new loop dimensions,
expanding the permutation term by orders of magnitude.
For the 11-layer fusion example validated in
Chapter~\ref{ch:validation}, the unified loop nest has 66
dimensions, yielding a map-space estimate exceeding $10^{515}$
(Section~\ref{sec:ps-mapspace}), compared with
${\sim}10^{24}\text{--}10^{32}$ for a single layer on the same
five-level, seven-total-level Eyeriss architecture.
The permutation term alone grows by over 400~orders of magnitude.

\paragraph{Extended Einsum notation.}
The use of the extended Einsum notation~\cite{een} as the workload
representation allowed the framework to capture fused convolutions,
including affine dimension coupling and stride-based cascading tile
dependencies, without workload-specific code paths.
This makes the tool applicable to any CONV-like layer topology
expressible in Einsum form.

\paragraph{Cost-model extensibility.}
The single-layer cost model (MOPs, energy via Accelergy, latency
with bandwidth-limited stall detection) was generalized to a
per-layer decomposition.
The resulting model correctly reproduces the DepFiN accelerator's
behavior: ideal latency of 37.94\,M cycles matches the
compute-bound reference exactly, DRAM bypass of intermediate
operands is enforced with zero intermediate DRAM traffic, and
on-chip memory stalls are negligible, arising only from the
boundary layers of the fusion pipeline (zero stalls at regime), as
expected.

\paragraph{Producer--consumer feasibility pruning.}
The pruning constraint, enforced at every hierarchy level through
the cascading tile-size inequality, eliminates infeasible mappings
before cost evaluation.
The pruning is exact: it removes only mappings that would violate
the producer--consumer data dependency, preserving every feasible
candidate.
Without such pruning, the fused-layer map space, exceeding
$10^{515}$ for the 11-layer configuration on Eyeriss, would be
entirely intractable.

% ====================================================================
%  §2  FUTURE WORK
% ====================================================================
\section{Future Work}
\label{sec:concl-future}

While the framework developed in this thesis demonstrates the
feasibility of analytical fused-layer mapping exploration, several
limitations remain that define clear directions for future research.

% ------------------------------------------------------------------
\subsection{Automated Mapper for Fused Workloads}
\label{sec:fw-mapper}

The most impactful limitation in the state of the art is the absence
of an automated mapper for multi-layer fused workloads
(Section~\ref{sec:ps-limitations}).
Currently, fused mappings must be specified manually through the
architecture definition file, which limits the exploration to
expert-guided configurations.
Extending the existing mapper infrastructure (linear, quadratic,
and hybrid search) to automatically explore the fused map space
will be feasible only with heuristics that reduce the search space,
even at the cost of potentially excluding some optimal solutions,
while respecting the producer--consumer pruning constraints at every
move.
This would transform the tool from an evaluator of prescribed
mappings into a full design-space exploration engine.

% ------------------------------------------------------------------
\subsection{Inter-Layer Parallelism}
\label{sec:fw-inter-layer}

The current model enforces sequential layer execution: all layers
share the same PE array and execute one at a time.
Architectures that exploit \emph{pipeline parallelism}, where
different layers compute concurrently on disjoint subsets of the
PE array, are not modeled.
The framework is well suited to this extension: supporting
inter-layer parallelism would require extending the latency model
to capture overlapping execution intervals, and would enable the
framework to evaluate a broader class of fused-layer schedules.

% ------------------------------------------------------------------
\subsection{Discovery of New Intra-Layer Reuse: Innermost Dimension-Sum Reuse}
\label{sec:fw-dimsum}

When two dimensions contribute to the same operand index through
an affine sum (e.g.\ $y + r$ in a convolution), the current model
does not exploit the overlap between consecutive iterations of the
sliding window within a single tile.
This leads to conservative (over-estimated) MOPs counts at the
innermost memory level.
Implementing fine-grained footprint analysis that accounts for
the sliding-window overlap would improve the accuracy of the
energy estimates, particularly for workloads with large filter
sizes.

% ------------------------------------------------------------------
\subsection{Porting to a Compiled Language}
\label{sec:fw-compiled}

Porting FactorFlow to a compiled systems programming language,
such as C++ or Rust, represents a highly promising avenue for
future work.
The current Python implementation, while effective for prototyping
and validation, incurs significant interpreter overhead and is
constrained by the Global Interpreter Lock (GIL), which prevents
true shared-memory multithreading.
Transitioning to a compiled environment would eliminate interpreter
overhead and enable lightweight, shared-memory multithreading over
process-based concurrency, significantly reducing execution time
and system overhead.
This improvement is particularly relevant for the feasibility of an
automated mapper.

% ------------------------------------------------------------------
\subsection{Integration with Optimus}
\label{sec:fw-graph-level}

The framework currently operates at the \emph{dataflow level},
optimizing the mapping of a fixed set of fused layers.
A complementary direction is \emph{graph-level} scheduling, which
decides \emph{which} layers to fuse and how to partition the
network graph into fusion segments.
Tools like Optimus~\cite{Optimus} address this partitioning
problem.
Coupling a graph-level scheduler with the fine-grained analytical
model developed in this thesis would enable joint optimization of
the fusion topology and the dataflow mapping, potentially
unlocking configurations that neither level can discover in
isolation.